% ===================================================================
%  TABELUL Nr. 3 — Copilăria și Tinerețea.
%  Traduction 2026 d'apres texto/io, controlee sur texto/fr, texto/en
%  et texto/es. Consigne : voir texto/ro/10-tabelul-01.tex.
%
%  Deux scenes, et l'ido subdivise beaucoup plus que le francais.
%
%  LE FAC-SIMILE OUBLIE DE METTRE « julio » EN GRAS au bloc c2-01-1.
%  La source ne bouge pas ; gravuri/korekti.json recoit sa ligne
%  roumaine — « (iulie sau ».
%
%  LES NOMS DES MOIS SONT LATINS ET NON SLAVES, contrairement a ce que
%  l'ukrainien fait a cote : « martie, aprilie, mai », « iunie, iulie,
%  august ». Le roumain a garde le calendrier romain la ou l'ukrainien
%  a garde le calendrier slave — deux voisins geographiques, deux
%  reponses opposees.
% ===================================================================

%%K t03-apar-1 apar
\VUsaut{3.0mm}
\VUtitre{15.0pt}{60}{TABELUL Nr. 3}
\VUsaut{1.6mm}
\VUtitre{11.4pt}{0}{Copilăria și Tinerețea.}
\VUsaut{3.4mm}

%%K t03-apar-2 apar
(Tabelele 3 și 4 sunt alcătuite astfel încât să dea, în patru \VUgras{scene de
familie}, fondul de cuvinte despre \VUgras{vreme}, \VUgras{vârstă},
\VUgras{familie}, \VUgras{îmbrăcăminte} și \VUgras{sănătate}.)
\VUsaut{4.0mm}

%%K t03-tit-1 sub
\VUcentre{12.0pt}{0}{\textit{Prima scenă.}}
\VUsaut{3.0mm}

%%K t03-tit-2 sub
\VUpk{10.2pt}{40}{Botezul la țară.}
\VUsaut{2.4mm}

%%K t03-tit-3 sub
\VUpk{10.2pt}{40}{Primăvara.}
\VUsaut{3.0mm}

%%K t03-c1-01-1 p
1. --- Suntem în \VUgras{primăvară}, în luna \VUgras{martie}, \VUgras{aprilie}
sau \VUgras{mai}, într-una din zilele \VUgras{săptămânii} --- \VUgras{luni},
\VUgras{marți} sau \VUgras{miercuri}. Este \VUgras{ora zece} dimineața.

%%K t03-c1-02-1 p
2. --- \VUgras{Vremea} este minunată, \VUgras{aerul} curat. Soarele
strălucește. Câțiva \VUgras{nori} mici \textsuperscript{(2)} se ridică pe
\VUgras{cerul} albastru \textsuperscript{(1)}.

%%K t03-c1-03-1 p
3. --- \VUgras{Copacii} nu au aproape deloc \VUgras{frunze}. \VUgras{Plopii}
zvelți \textsuperscript{(3)} și \VUgras{platanul} înalt \textsuperscript{(17)}
sunt acoperiți deocamdată numai cu muguri.

%%K t03-c1-04-1 p
4. --- \VUgras{Rândunelele} \textsuperscript{(4)} s-au întors și se rotesc cu
strigăte vesele în jurul \VUgras{săgeții} \textsuperscript{(7)}
\VUgras{clopotniței} \textsuperscript{(5)}, care încununează \VUgras{biserica}
satului \textsuperscript{(6)}.

%%K t03-tit-4 sub
\VUsaut{3.4mm}
\VUpk{10.2pt}{40}{Biserica.}
\VUsaut{3.0mm}

%%K t03-c1-05-1 p
5. --- Foarte simplă este această biserică, cu \VUgras{portalul} ei mare
\textsuperscript{(13)} și cu \VUgras{ceasul} ei mare \textsuperscript{(8)}.
\VUgras{Limba mică} \textsuperscript{(9)} stă pe cifra X: ea arată
\VUgras{orele}. \VUgras{Cea mare} \textsuperscript{(10)} stă pe XII (amiaza); ea
arată \VUgras{minutele}.

%%K t03-c1-06-1 p
6. --- În clopotniță sună vesel \VUgras{clopotele} \textsuperscript{(11)}. În
vârful acoperișului stă o \VUgras{cruce} mică de piatră \textsuperscript{(12)}.

%%K t03-c1-07-1 p
7. --- Biserica nu are numai portalul cel mare \textsuperscript{(13)}; are și o
ușă mică laterală. Prin această ușă \VUgras{preotul} \textsuperscript{(15)}, în
\VUgras{sutana} sa, iese din \VUgras{sacristie} \textsuperscript{(14)} și se
îndreaptă spre \VUgras{casa preotului} \textsuperscript{(19)}.

%%K t03-c1-08-1 p
8. --- Lângă el se află \VUgras{lăptăreasa} \textsuperscript{(24)} cu
\VUgras{măgarul} ei \textsuperscript{(23)}. Se întoarce din oraș, unde a dus
copiilor laptele ei bun.

%%K t03-tit-5 sub
\VUsaut{4.0mm}
\VUpk{10.2pt}{40}{Livada.}
\VUsaut{3.4mm}

%%K t03-c1-09-1 p
9. --- Domnul Sur (măgarul) este foarte mulțumit de această plimbare de dimineață.
Trage cu desfătare aerul înmiresmat de \VUgras{floarea} \textsuperscript{(20)}
care acoperă \VUgras{pomii} \textsuperscript{(18)} din livada vecină. Toți în roz
și alb sunt acești \VUgras{piersici} \textsuperscript{(18)} și acești
\VUgras{caiși} \textsuperscript{(19)}, și făgăduiesc o recoltă bună.

%%K t03-c1-10-1 p
10. --- Iar în vremea aceasta \VUgras{albinele} mici \textsuperscript{(22)} din
\VUgras{stup} \textsuperscript{(21)} zboară acolo să-și strângă \VUgras{mierea}
dulce, bâzâind la lucru.

%%K t03-c1-11-1 p
11. --- Dar ce s-a întâmplat? Iată că vin în fugă copiii satului și împrăștie
\VUgras{gâștele} speriate \textsuperscript{(25)}. Acest alai iese din biserică
după \VUgras{botezul} fiului notarului.

%%K t03-c1-12-1 p
12. --- Micul Iacob se trezește în \VUgras{leagănul} său \textsuperscript{(26)}
din sunetul vesel al clopotelor, iar sora lui Magdalena, care vrea și ea să vadă
alaiul, o roagă pe mama să o pieptene și să o îmbrace în prag.

%%K t03-c1-13-1 p
13. --- Mama se învoiește. Își pune \VUgras{basmaua} \textsuperscript{(28)},
\VUgras{bluza} \textsuperscript{(29)} și \VUgras{șorțul} \textsuperscript{(30)}
și se așază pe un scaun scund înaintea ușii. Magdalena are pe ea numai
\VUgras{fusta} \textsuperscript{(31)} și \VUgras{corsajul}
\textsuperscript{(32)}.

%%K t03-c1-14-1 p
14. --- Ce fericită este! Își uită florile cele mai dragi, \VUgras{garoafele}
\textsuperscript{(34)}, pe care se așază \VUgras{fluturii}
\textsuperscript{(35)}, \VUgras{lalelele} \textsuperscript{(36)},
\VUgras{lăcrămioarele} \textsuperscript{(37)} din \VUgras{ghivece}
\textsuperscript{(38)} de pe \VUgras{zid} \textsuperscript{(33)} și
\VUgras{trandafirii} ei \textsuperscript{(39)}, care alcătuiesc un gard atât de
frumos. Se uită lung la veșmintele bogate ale doamnelor din alai.

%%K t03-c1-15-1 p
15. --- Băieții satului se uită puțin la veșminte. Se reped înainte și se
îmbrâncesc ca să apuce \VUgras{bomboane} \textsuperscript{(55)} sau
\VUgras{monede} \textsuperscript{(52)}. Unul dintre ei plânge
\textsuperscript{(40)}.

%%K t03-c1-16-1 p
16. --- Altul a călcat pe laba unui \VUgras{câine} \textsuperscript{(42)}, iar
acesta fuge schelălăind și pune pe goană o \VUgras{găină}
\textsuperscript{(43)} cu \VUgras{puii} ei \textsuperscript{(44)}.

%%K t03-tit-6 sub
\VUsaut{5.0mm}
\VUpk{10.2pt}{40}{Alaiul botezului.}
\VUsaut{4.0mm}

%%K t03-c1-17-1 p
17. --- În fruntea alaiului merge \VUgras{doica} \textsuperscript{(45)}. Poartă o
\VUgras{scufie} frumoasă \textsuperscript{(46)} cu panglici lungi. Îl duce pe
\VUgras{nou-născut} \textsuperscript{(47)}, tot în \VUgras{dantele}
\textsuperscript{(48)}.

%%K t03-c1-18-1 p
18. --- În urma doicii vin \VUgras{nașul} \textsuperscript{(49)} și
\VUgras{nașa} \textsuperscript{(53)}. Împart bomboane și monede. Nașul are
\VUgras{mănuși} \textsuperscript{(51)} și \VUgras{joben} \textsuperscript{(50)}.
Nașa ține în mână un \VUgras{buchet} frumos \textsuperscript{(54)}.

%%K t03-c1-19-1 p
19. --- În urma lor se văd \VUgras{unchiul} \textsuperscript{(56)} și
\VUgras{mătușa} \textsuperscript{(58)} copilului. Mătușa are o \VUgras{pălărie}
elegantă \textsuperscript{(59)}, unchiul o \VUgras{redingotă} neagră
\textsuperscript{(57)} și o vestă albă. Mai departe vin \VUgras{vărul}
\textsuperscript{(60)} și \VUgras{verișoara} \textsuperscript{(61)}. Aceasta din
urmă duce de mână pe \VUgras{sora} \textsuperscript{(62)} nou-născutului, pe mica
Berta.

%%K t03-c1-20-1 p
20. --- Celălalt \VUgras{frate} \textsuperscript{(63)} al Bertei merge în urmă. Dă
brațul unei \VUgras{rude} \textsuperscript{(64)}. În urma lor vin
\VUgras{prietenii} \textsuperscript{(65)} familiei.

%%K t03-tit-7 sub
\VUsaut{4.6mm}
\VUpk{10.2pt}{40}{Gură-cască.}
\VUsaut{3.4mm}

%%K t03-c1-21-1 p
21. --- Lângă alai se află un \VUgras{cerșetor} \textsuperscript{(66)}, sprijinit
în \VUgras{cârja} sa \textsuperscript{(67)}. Cere de pomană.

%%K t03-c1-22-1 p
22. --- De cealaltă parte se află un băiat foarte \VUgras{politicos}
\textsuperscript{(68)}. Își ridică șapca și le spune \VUgras{«bună ziua»} celor pe
care îi cunoaște.

%%K t03-c1-23-1 p
23. --- Sătenii au ieșit din case ca să vadă alaiul botezului. Vedem un
\VUgras{țăran} \textsuperscript{(69)} cu \VUgras{tichie de hârtie}, iar lângă el
o \VUgras{țărancă} \textsuperscript{(70)}. Apoi o \VUgras{bătrână}
\textsuperscript{(71)}, sprijinită în \VUgras{bastonul} ei \textsuperscript{(72)}.
În sfârșit \VUgras{morarul} \textsuperscript{(73)} cu o \VUgras{pălărie de
pâslă} largă \textsuperscript{(74)} și o țărancă tânără în \VUgras{saboți}
\textsuperscript{(75)}, care își învață copilașul să meargă.

%%K t03-c1-24-1 p
24. --- Tabloul este foarte însuflețit.
\VUsaut{4.0mm}

%%K t03-tit-8 sub
\VUcentre{12.0pt}{0}{\textit{A doua scenă.}}
\VUsaut{3.0mm}

%%K t03-tit-9 sub
\VUpk{10.2pt}{40}{Sărbătoarea populară.}
\VUsaut{2.4mm}

%%K t03-tit-10 sub
\VUpk{10.2pt}{40}{Întreceri pe apă și regate.}
\VUsaut{3.0mm}

%%K t03-c2-01-1 p
1. --- A venit \VUgras{vara}. Soarele arzător al lunii \VUgras{iunie} (iulie sau
\VUgras{august}) îi îmbie pe tineri să se scalde și să se plimbe cu barca.

%%K t03-c2-02-1 p
2. --- Pe malul drept al râului vâslașii se dezbracă, ca să ia parte la întrecerile
pe apă și la regate.

%%K t03-c2-03-1 p
3. --- Unul dintre ei își scoate \VUgras{ghetele} \textsuperscript{(1)}; le
descheie (le desface). Doi dintre tovarășii lui sunt gata. Acum au pe ei numai
\VUgras{maioul} \textsuperscript{(2)} și \VUgras{slipul} \textsuperscript{(3)}.
Stau de vorbă, așteptându-și tovarășii. Vorbesc despre bâlci și despre serbarea de
pe celălalt mal.

%%K t03-c2-04-1 p
4. --- Pe pământ lângă ei stau \VUgras{hainele} tovarășilor lor: o
\VUgras{pereche} de \VUgras{cizme} \textsuperscript{(4)}, \VUgras{papuci}
\textsuperscript{(5)}, \VUgras{ciorapi} \textsuperscript{(6)}, pălăriile de
\VUgras{pâslă} \textsuperscript{(7)} și de \VUgras{paie} \textsuperscript{(8)} și
o \VUgras{cravată} \textsuperscript{(9)}.

%%K t03-c2-05-1 p
5. --- Unul dintre tineri și-a împăturit cu grijă \VUgras{haina}
\textsuperscript{(10)}. O pune pe un ștergar, în care vecinul său va înfășura
îndată amândouă \VUgras{costumele} lor.

%%K t03-c2-06-1 p
6. --- Al șaselea băiat întârzie. Pe cap are încă \VUgras{șapca}
\textsuperscript{(11)}. Nu și-a scos nici \VUgras{cămașa} \textsuperscript{(12)},
nici \VUgras{gulerul} \textsuperscript{(13)}, nici \VUgras{manșetele}
\textsuperscript{(14)}. Ține în mână \VUgras{vesta} \textsuperscript{(17)} și este
încă în \VUgras{pantaloni} \textsuperscript{(16)} și în \VUgras{bretele}
\textsuperscript{(15)}. La cursa următoare nu va ajunge, de bună seamă.

%%K t03-c2-07-1 p
7. --- Lângă tineri o potecă năpădită de \VUgras{ciulini} \textsuperscript{(18)}
coboară spre apă, unde bătrânul Iacob pregătește printre \VUgras{stuf}
\textsuperscript{(19)} \VUgras{focurile de artificii} \textsuperscript{(20)}, care
vor fi slobozite seara.

%%K t03-tit-11 sub
\VUsaut{4.4mm}
\VUpk{10.2pt}{40}{Regata.}
\VUsaut{3.4mm}

%%K t03-c2-08-1 p
8. --- Pe râu câteva doamne, adăpostite sub umbrele, se plimbă cu \VUgras{barca}
\textsuperscript{(21)}. Urmăresc regata, iar aceasta este foarte plăcută. În
fiecare \VUgras{șalupă} \textsuperscript{(22)} patru \VUgras{vâslași}
\textsuperscript{(24)} vâslesc din răsputeri, iar \VUgras{cârmaciul}
\textsuperscript{(26)} ține \VUgras{cârma} \textsuperscript{(25)} și conduce
gingașa ambarcațiune. \VUgras{Vâslele} \textsuperscript{(23)} bat apa în cadență,
și bărcile ușoare zboară cu o iuțeală amețitoare.

%%K t03-c2-09-1 p
9. --- Câțiva tineri se întrec lângă piciorul podului pentru premiul
\VUgras{catargului cu premiu} \textsuperscript{(28)}. Unul dintre ei înaintează cu
grijă pe prăjina mlădioasă și alunecoasă, ca să apuce \VUgras{stegulețul}
\textsuperscript{(29)} prins la capăt. \VUgras{Potrivnicul} său nenorocos
\textsuperscript{(30)} a căzut în apă. Se întoarce înot la mal.

%%K t03-c2-10-1 p
10. --- Băieții mai mari se \VUgras{luptă pe apă} \textsuperscript{(32)} lângă
\VUgras{cheiul} \textsuperscript{(33)}. O \VUgras{mulțime} mare
\textsuperscript{(31)} privește toate aceste jocuri, sprijinită pe
\VUgras{balustrada} \VUgras{podului} \textsuperscript{(27)} și înghesuindu-se
de-a lungul \VUgras{malului}. Unii gură-cască se întorc totuși ca să privească
jocul cu \VUgras{stâlpul} \textsuperscript{(45)} sau ca să admire \VUgras{alaiul
primăriei}, care merge la \VUgras{împărțirea premiilor}.

%%K t03-c2-11-1 p
11. --- Mai întâi aleargă câțiva \VUgras{băiețandri} \textsuperscript{(35)}
înaintea \VUgras{muzicanților} \textsuperscript{(36)}; apoi vin \VUgras{primarul}
\textsuperscript{(37)}, ajutoarele lui și \VUgras{consiliul comunal} al
târgului. La urmă pășesc \VUgras{pompierii} \textsuperscript{(38)}.

%%K t03-tit-12 sub
\VUsaut{5.0mm}
\VUpk{10.2pt}{40}{Bâlciul.}
\VUsaut{3.6mm}

%%K t03-c2-12-1 p
12. --- În \VUgras{piața bâlciului} \textsuperscript{(39)} este mare înghesuială.
Lumea vine să vadă felurite \VUgras{barăci}: \VUgras{călușeii}
\textsuperscript{(40)}, \VUgras{circul} \textsuperscript{(41)}, unde
reprezentația a și început; \VUgras{balul popular} \textsuperscript{(42)}, unde
tinerii și fetele din târg se dau plăcerii de a \VUgras{dansa}.

%%K t03-c2-13-1 p
13. --- În fund se vede \VUgras{arena} \textsuperscript{(44)}, unde îndrăznețul
\VUgras{matador} spaniol se luptă cu \VUgras{taurul} înfuriat
\textsuperscript{(45)}; apoi \VUgras{munții ruși} \textsuperscript{(49)} și
\VUgras{tirul} \textsuperscript{(46)}, unde se deprind cu \VUgras{armele de foc}
și cu tragerea la \VUgras{țintă}.

%%K t03-c2-14-1 p
14. --- Alături se află \VUgras{teatrul de bâlci} \textsuperscript{(47)}, unde se
joacă \VUgras{comedie} și \VUgras{dramă}. Foarte aproape stă baraca
\VUgras{uriașului} \textsuperscript{(48)} și a \VUgras{piticului}.
\VUgras{Statura} celui dintâi este de doi metri și cincisprezece; al doilea este
abia cât un copil.

%%K t03-c2-15-1 p
15. --- În sfârșit, iată marea \VUgras{menajerie} \textsuperscript{(50)} cu
\VUgras{fiarele} ei sălbatice, cu \VUgras{tigrul} \textsuperscript{(51)} și cu
\VUgras{ghearele} lui puternice, cu \VUgras{leul} \textsuperscript{(52)} și cu
\VUgras{coama} lui deasă, cu doi \VUgras{girafi} \textsuperscript{(53)} și cu
\VUgras{elefantul} \textsuperscript{(54)}, care își mânuiește atât de îndemânatic
\VUgras{trompa}.

%%K t03-c2-16-1 p
16. --- \VUgras{Îmblânzitorul} \textsuperscript{(55)}, care dresează aceste fiare,
stă la ușa barăcii.
\VUsaut{6.0mm}
\VUfilet{22mm}
